\section{Setting}
\label{sec:setting}

\paragraph{Models.} Holo-3.1 models are fine-tunes of the Qwen family \citep{hcompany2026holo31}; their released configuration is a Qwen3.5 vision-language model. Their decoder is hybrid: three linear-attention layers (Gated DeltaNet~\citep{yang2025gateddeltanet}) for every full-attention layer, and tied input and output embeddings. Holo-3.1-4B has \cfgBigLayers{} layers of width \cfgBigHidden{} (MLP width \cfgBigInter), \cfgBigNLinear{} of them Gated DeltaNet and \cfgBigNFull{} full attention; Holo-3.1-0.8B has \cfgSmallLayers{} layers of width \cfgSmallHidden. Layers are indexed from zero. At each step the agent writes one JSON object whose fields come in a fixed order, a \texttt{note}, a \texttt{thought}, then the \texttt{tool\_call}, after its reasoning. % source: data/pod_hf/config_holo31_4B.json and config_holo31_0.8B.json (figs/b4.py); ~/osw/rtg/Agent/holo_agent/tools.py make_step_model (note, thought, tool_call)

\paragraph{OSWorld harness (4B).} We reuse the Holo agent loop as packaged by RealtimeGUI at commit \texttt{a763361} \citep{realtimegui}, a third-party reproduction rather than an H Company repository, changing only the model provider. Every run uses its default system prompt (variant \texttt{v2}, recorded in each episode's metadata), which is RealtimeGUI's, not H Company's, and already asks the agent to check each action on the new screenshot, to ``never assume success'' and to call \texttt{answer} only once it has ``verified the result on screen''. False ``done'' is therefore measured under a prompt that already asks for verification. % source: ~/osw/rtg/Agent/holo_agent/agent.py _SYSTEM_PROMPT ("v2", the default); data/local/harness/run_holo.py:248 (--prompt default v2); prompt_variant "v2" in all 120 meta.json of results_test_gh_orig_modal_test, _orig_guard_claude, _merge_cua_b050
Its rule for recognising a declared infeasibility (\texttt{task.\{0,4\}infeasible|infeasible|cannot be (completed|done)|not possible}) is broader than that of the official OSWorld adapter for Holo, which accepts only \texttt{task.\{0,4\}infeasible}; our absolute scores are therefore not strictly comparable with H Company's, while paired comparisons within our harness are. % source: data/local/harness/run_holo.py (_INFEASIBLE_RE), data/local/harness/surfer_agent.py (official xlang-ai/OSWorld mm_agents/surferH); git remote of ~/osw/rtg = github.com/1azybug/RealtimeGUI @a763361; data/local/JOURNAL.md 02:25 UTC
vLLM~0.23 serves the 4B with thinking enabled, $T=0.8$, at most 50 steps and the three latest screenshots in context. % source: JOURNAL.md "OSWorld — infra prête" and "OSWorld — pilote"
OSWorld \citep{xie2024osworld} tasks run in Ubuntu VMs under Docker+KVM, at most four in parallel, and are scored by the benchmark's execution-based checkers. We fixed three infrastructure issues in the VM provider (a memory cap, binding VM ports to localhost, an nftables network script) and modified no task or checker ourselves. The OSWorld copy inside RealtimeGUI, which every local and GitHub run used, is not byte-identical to upstream: its Chrome evaluator getters wait for page load with a fallback (\texttt{networkidle}, then \texttt{load}, then \texttt{domcontentloaded}) instead of \texttt{networkidle} alone, a change its authors made for a slow proxy. Checkers of Chrome tasks that fetch a web page may therefore score differently from upstream OSWorld. % source: JOURNAL.md "3 correctifs d'infra du provider docker"; ~/osw/rtg @a763361 Env/OSWorld/desktop_env/evaluators/getters/chrome.py (_WAIT_PLAN, _goto_with_fallback; commit c1feb74), ~/osw/rtg/复现_debug_评测器修复.md; holo-osworld-runner .github/workflows/osworld.yml (RealtimeGUI @a763361 + patches/rtg-a763361.patch, which touches only provider.py and run_holo.py)

\paragraph{Frozen split.} From the 361 tasks of RealtimeGUI's \texttt{test\_nogdrive.json}, we sort ids by SHA-1 within each domain and deal them alternately into a \textsc{diag} and a \textsc{test} pool, keeping 120 tasks each with per-domain quotas. The split files have SHA-1 prefixes \texttt{51e670b3} (diag) and \texttt{957d68e7} (test). No construction or screening read the test split. Six valid online runs then used it: five model runs (\S\ref{sec:online}), the original, served once by our H200 server and once by Modal, the ExPO edit, which went to it without having passed the gate, a partial corrective-prompt run, and the CUA-Lite merge under a decision rule fixed after its reserve scores had been seen; and the guarded agent (\S\ref{sec:guard}). A seventh run, of an output-head scaling, was invalidated (\S\ref{sec:online}). % source: JOURNAL.md "OSWorld — split figé"; data/local/JOURNAL.md 03:05 UTC; data/local/reports/REDTEAM_0305.md point 5

\paragraph{MiniWoB++ harness (0.8B).} Our first attempt used Holo-3.1-0.8B on MiniWoB++ \citep{shi2017wob,liu2018miniwob} through BrowserGym \citep{browsergym2024}, with the task area cropped and upscaled $5\times$; it is described in \S\ref{sec:mw} and Appendix~\ref{app:mw}.

\paragraph{Hardware and licences.} Rollouts, replays and screens ran on one NVIDIA H200 shared between the vLLM server and research jobs; edited checkpoints were also served on rented H100 GPUs, and the OSWorld VMs ran on a separate CPU host. As order-of-magnitude estimates, the project used about 14 hours of one H200 (rented, shared by every research job and the rollout server, so per-experiment GPU-hours are not separable), about 8 H100 GPU-hours (USD 31.9 billed over two Modal accounts by 05:11 UTC) and about 60 GitHub Actions CPU jobs for the OSWorld VMs. These figures stop at 05:11 UTC: the later Modal deployments (the Holo4-27B judge and agent, the Holo-3.1 model behind the guarded agent) and the Claude calls of the verifier are not included. % source: data/local/MODAL_COST.md (16.65 + 15.26 USD at 05:11 UTC, H100 at about 4 USD/h); data/local/modal_deploys/20261108T075938_acct2_Holo4-27B.json, 20261108T090300_acct2_Holo4-27B-agent.json (after 05:11); H200 hours and job count: coordinator estimates (ares, 05:25 UTC)
% source: data/local/JOURNAL.md ("vLLM 0.23 sur H200", 01:10 UTC "Modal H100"); data/pod_astra/merge/VERIFY_CUA_B050_BEHAVIOR.md (H200 vs H100)
Holo-3.1-0.8B and -4B, Qwen3.5-4B, the CUA-Lite donor, OSWorld and BrowserGym are released under Apache-2.0 and MiniWoB++ under MIT; the RealtimeGUI repository has no licence file at the cited commit, and we use it unmodified except for the model provider. % source: Hugging Face model API and GitHub license API, queried 2026-11-08
